\scanpage{13}
\uline{Remarques} a) Si $T$ peut être tordu, on ne peut plus affirmer que l'hom injectif
\[
\mathrm{Pic}(D)\longrightarrow\mathrm{Pic}(\mathscr B)
\]
soit surjectif. Mais (si $S$ \deleted{\unsure{est noeth.})} n'a qu'un nb fini de composantes connexes) on voit que le \uline{conoyau est fini} en fait on a \deleted{pour} $S$ étant de \unsure{dbl|ctible} une suite exacte
\[
\begin{aligned}
0-\mathrm{Pic}(D)-\mathrm{Pic}(\mathscr B)&\longrightarrow H^1(\Gamma,M_{\bar{\unsure{\xi|\zeta}}}),\\[-1pt]
&\hspace{8mm}\Vert\\[-2pt]
&\mathrm{Pic}(\mathscr B_\eta)\xrightarrow{\sim}\mathrm{Pic}(T_\eta)\xleftarrow{\sim}\mathrm{Pic}(T_\xi)
\end{aligned}
\]
où $\eta$ pt gén. de $D$\quad $\xi$ pt gén. de $S$

\deleted{\illegible{1}} où $\Gamma$ est l'image de $\pi_1(S,\bar\xi)$ dans $\operatorname{Aut}(M_{\bar\xi})$.

Si $S$ est régulier, \deleted{\unsure{et alors}} on peut mettre un zéro à droite, i.e. la condition $H^1(\Gamma,M_{\bar\xi})=0$ est néc. et suff. pour que l'on ait bijectivité $\mathrm{Pic}(D)-\mathrm{Pic}(\mathscr B)$ \deleted{\unsure{et}} \deleted{\unsure{elle implique en tous cas inject.}}\inserted{ou même} bijectivité $\mathrm{Pic}(D)\simeq\mathrm{Pic}(\mathscr B^{\mathrm{reg}})$ [car $\mathrm{Pic}(\mathscr B_\eta)\longrightarrow\mathrm{Pic}(\mathscr B_{\unsure{p|\eta|\xi}}^{\mathrm{reg}})$ est \inserted{en tous cas} \unsure{bijectif|injectif|surjectif} pour des raisons de codim $\geq 2$ et de régularité]

b) L'injectivité de $\mathrm{Pic}(D)-\mathrm{Pic}(\mathscr B)$ et la caractérisation des éléments de $\mathrm{Pic}(D)$ amples/$S$ comme ceux dont l'image inverse sur $\mathscr B$ sont amples /$S$, restent vrais sans hyp. de normalité sur $S$.

c) Supposons $S$ normal \inserted{et $H^1(\xi,M)=0$ p.ex. $T$ non tordu}. Alors
\[
\mathrm{Pic}(S)\xrightarrow{\sim}\mathrm{Pic}(G')
\]

En effet, on \deleted{sait que $\mathrm{Pic}(G')$} a injectivité car $G'$ a une section sur $S$ (section unité). Reste à prouver que si $\unsure{\mathscr F|\varphi|\xi}\in\mathrm{Pic}(G')$ induit 0 sur la section unité, alors $\mathscr F$ est nul. Or
\[
\mathrm{Pic}(G')\hookrightarrow\mathrm{Pic}(G'^{\mathrm{reg}})\bigl[\simeq\mathrm{Pic}(\mathscr B^{\unsure{\mathrm{reg}|\prime\mathrm{reg}}})\xleftarrow{\sim}\mathrm{Pic}(\mathscr B)\bigr]
\]
inj. pour raisons de profondeur, donc il suffit \deleted{que} de voir que l'image \deleted{\unsure{inv|inverse}} de $\mathscr F$ dans $\mathrm{Pic}(\mathscr B)$ est nulle. Mais il suffit de voir que la restriction de cette dernière à la section unité de $\mathscr B$ sur $D$ est nulle. Mais $D\longrightarrow G'$
% archive sheet number: 12
